%%%%%%%%%%%%%%%%%%%%%%%%%%%%%%%%%%%%%%%%%%%%%%%%%
%%%%%%%%%%%% chap: results %%%%%%%%%%%%%%%%%%%%
%%%%%%%%%%%%%%%%%%%%%%%%%%%%%%%%%%%%%%%%%%%%%%%%%

\chapter{Experimental Results and Evaluation}\label{chapter:chap4}

This chapter presents the results of deploying the multi-stage phishing
honeypot and running the machine learning pipeline on the collected dataset.
The evaluation addresses three questions: whether the system successfully
captures the full attacker lifecycle, whether the extracted behavioural
features distinguish attacker profiles, and how well the Random Forest
classifiers perform on both classification tasks.

\section{Data Collection}

The system ran continuously for approximately 54 hours with all four bot
profiles active. Table~\ref{tab:dataset_summary} summarises the resulting
dataset. A total of 38,438 session records were written to the
\texttt{ml\_features} table, split evenly between credential submission
and SSH attack stages, confirming that the two-phase attack model
executed successfully for every cycle.

\begin{table}[h]
\centering
\caption{Dataset summary by bot type and attack stage.}
\label{tab:dataset_summary}
\begin{tabular}{lrrr}
\hline
\textbf{Bot type} & \textbf{Cred. submission} & \textbf{SSH attack} & \textbf{Total} \\
\hline
Scanner            & 8,848 & 8,847 & 17,695 \\
Credential stuffer & 9,210 & 9,208 & 18,418 \\
Human-like         & 522   & 519   & 1,041  \\
Sophisticated      & 644   & 640   & 1,284  \\
\hline
\textbf{Total}     & 19,224 & 19,214 & 38,438 \\
\hline
\end{tabular}
\end{table}

The class imbalance between scanner/credential stuffer and the slower
human-like/sophisticated profiles is a direct consequence of the bot cycle
timings: the scanner executes a full cycle every 15 seconds while the
sophisticated profile, with its deliberate inter-command delays, requires
several minutes per cycle. This distribution reflects the real-world
prevalence of automated scanning relative to targeted APT activity
\cite{Bringer2012HoneypotSurvey}, and is handled in training via
\texttt{class\_weight=balanced}.

\section{Behavioural Feature Analysis}

Before examining classifier performance, it is instructive to inspect
the raw feature distributions, since these form the basis of all
downstream classification. Table~\ref{tab:feature_means} shows the
mean value of each feature per bot type.

\begin{table}[h]
\centering
\caption{Mean feature values per bot type across all collected sessions.}
\label{tab:feature_means}
\begin{tabular}{lrrrrrr}
\hline
\textbf{Profile} & \textbf{time\_on\_page} & \textbf{typing\_delay} & \textbf{field\_corr.} & \textbf{cmd\_count} & \textbf{duration\_s} & \textbf{cmd\_delay\_s} \\
 & \textit{(ms)} & \textit{(ms)} & & & & \\
\hline
Scanner            & 149    & 25    & 0.00 & 1.49 & 1.95   & 0.14 \\
Cred. stuffer      & 630    & 50    & 0.00 & 1.43 & 2.40   & 0.32 \\
Human-like         & 3,843  & 749   & 0.52 & 3.15 & 18.08  & 2.41 \\
Sophisticated      & 8,941  & 1,586 & 0.58 & 7.73 & 101.11 & 6.03 \\
\hline
\end{tabular}
\end{table}

Several observations are noteworthy. First, the four profiles are
\textbf{monotonically ordered} across all six features: every feature
value increases from scanner through to sophisticated. This ordering is
consistent with the CyberLab distributions used to parameterise the
bots \cite{Alata2006Honeynet}, providing a validation that the simulation
faithfully reproduces real-world attacker timing patterns.

Second, the separation between profiles is \textbf{large in absolute
terms}. The sophisticated profile spends on average 8.9 seconds on the
phishing page versus 149 milliseconds for the scanner — a factor of
approximately 60. SSH session duration shows an even wider spread:
101.11 seconds for sophisticated versus 1.95 seconds for scanner, a
factor of 52. These large margins explain the high classification
accuracy achieved without requiring complex feature engineering.

Third, \texttt{field\_corrections} is non-zero only for human-like and
sophisticated profiles, reflecting the realistic typing simulation
implemented for those bots. Scanner and credential stuffer submit forms
programmatically with no corrections, which is consistent with their
tool-based nature.

\section{Classification Results}

\subsection{RF1 — Bot Type Classification}

The first Random Forest classifier predicts which of the four attacker
profiles generated a given session. Table~\ref{tab:rf1_cv} reports
5-fold stratified cross-validation results on the full 38,438-row dataset.

\begin{table}[h]
\centering
\caption{RF1 (bot type) — 5-fold stratified cross-validation results.}
\label{tab:rf1_cv}
\begin{tabular}{lrr}
\hline
\textbf{Metric} & \textbf{Mean} & \textbf{Std} \\
\hline
Accuracy    & 97.87\% & $\pm$0.11\% \\
F1 weighted & 97.87\% & $\pm$0.11\% \\
F1 macro    & 98.83\% & $\pm$0.08\% \\
\hline
\end{tabular}
\end{table}

The macro F1 of 98.8\% indicates that the classifier performs
consistently well across all four classes, including the minority
human-like (1,041 rows) and sophisticated (1,284 rows) profiles. The standard
deviation of $\pm$0.1\% confirms that this performance is stable across
folds.

Table~\ref{tab:rf1_report} shows the per-class precision, recall, and
F1 score evaluated on the full training set.

\begin{table}[h]
\centering
\caption{RF1 (bot type) — per-class classification report (full dataset).}
\label{tab:rf1_report}
\begin{tabular}{lrrrr}
\hline
\textbf{Class} & \textbf{Precision} & \textbf{Recall} & \textbf{F1} & \textbf{Support} \\
\hline
Scanner            & 98.7\% & 99.3\% & 99.0\% & 17,695 \\
Credential stuffer & 99.3\% & 98.7\% & 99.0\% & 18,418 \\
Human-like         & 100.0\% & 99.9\% & 100.0\% & 1,041  \\
Sophisticated      & 100.0\% & 100.0\% & 100.0\% & 1,284  \\
\hline
\textbf{Macro avg} & 99.5\% & 99.5\% & 99.5\% & 38,438 \\
\hline
\end{tabular}
\end{table}

The sophisticated profile achieves perfect precision and recall. This
result is consistent with the feature analysis: the sophisticated bot's
timing parameters (mean page time of 8.9 s, mean inter-command delay of
6.0 s) are sufficiently distant from all other profiles that no
misclassification occurs. The small number of errors in the scanner and
credential stuffer classes reflect occasional overlap at the boundary
of their timing distributions.

\subsection{RF2 — Attack Stage Classification}

The second classifier predicts whether a session belongs to the
credential submission or SSH attack stage.
Table~\ref{tab:rf2_cv} shows perfect cross-validation performance
across all metrics.

\begin{table}[h]
\centering
\caption{RF2 (attack stage) — 5-fold stratified cross-validation results.}
\label{tab:rf2_cv}
\begin{tabular}{lrr}
\hline
\textbf{Metric} & \textbf{Mean} & \textbf{Std} \\
\hline
Accuracy    & 100.00\% & $\pm$0.00\% \\
F1 weighted & 100.00\% & $\pm$0.00\% \\
F1 macro    & 100.00\% & $\pm$0.00\% \\
\hline
\end{tabular}
\end{table}

This result warrants explanation. The two attack stages are linearly
separable by design: credential submission records have zero values for
all SSH features (\texttt{command\_count}, \texttt{session\_duration\_s},
\texttt{inter\_cmd\_delay\_avg\_s}), while SSH attack records have zero
values for all phishing features (\texttt{time\_on\_page\_ms},
\texttt{typing\_delay\_ms}, \texttt{field\_corrections}). The classifier
is therefore exploiting a structural property of the dataset encoding
rather than learning a generalised decision boundary. This is not a
flaw in the methodology — the two stages represent genuinely distinct
event types at different points in the attack lifecycle — but it means
the RF2 result should be interpreted as a validation of the schema
design rather than a measure of classifier generalisation.

\section{Feature Importance Analysis}

Figure~\ref{fig:feature_importance} (see \texttt{results/feature\_importance.png})
shows the Gini importance of each feature in both classifiers.
Table~\ref{tab:feature_importance} summarises the values.

\begin{table}[h]
\centering
\caption{Gini feature importances for both Random Forest classifiers.}
\label{tab:feature_importance}
\begin{tabular}{lrr}
\hline
\textbf{Feature} & \textbf{RF1 (bot type)} & \textbf{RF2 (attack stage)} \\
\hline
\texttt{inter\_cmd\_delay\_avg\_s} & 25.2\% & 15.7\% \\
\texttt{typing\_delay\_ms}         & 24.6\% & 30.0\% \\
\texttt{time\_on\_page\_ms}        & 20.1\% & 30.5\% \\
\texttt{session\_duration\_s}      & 17.2\% & 15.7\% \\
\texttt{command\_count}            & 8.3\%  & 8.1\%  \\
\texttt{field\_corrections}        & 4.7\%  & 0.0\%  \\
\hline
\end{tabular}
\end{table}

For RF1, the three most important features — \texttt{inter\_cmd\_delay\_avg\_s}
(25.2\%), \texttt{typing\_delay\_ms} (24.6\%), and
\texttt{time\_on\_page\_ms} (20.1\%) — together account for nearly 70\%
of the classification signal. Crucially, no single feature dominates;
the importance is distributed across both phishing-stage and SSH-stage
features. This indicates that the classifier is genuinely using the
cross-stage behavioural fingerprint rather than relying on a single
easily-spoofed indicator \cite{Holz2006AttackPatterns}.

The absence of \texttt{field\_corrections} importance in RF2 (0.0\%)
is expected: both attack stages contain sessions from all four profiles,
so this within-profile feature provides no discriminating information
for stage classification.

\section{Comparison with CICIDS2017 Baseline}

To contextualise the results, Table~\ref{tab:baseline} compares the
honeypot RF1 performance against published Random Forest results on the
CICIDS2017 benchmark \cite{Bringer2012HoneypotSurvey}, the standard
reference dataset for network intrusion classification.

\begin{table}[h]
\centering
\caption{Comparison of honeypot RF1 results against CICIDS2017 baseline.}
\label{tab:baseline}
\begin{tabular}{llrr}
\hline
\textbf{System} & \textbf{Task} & \textbf{Accuracy} & \textbf{F1 macro} \\
\hline
CICIDS2017 (RF, published range) & 14-class network flow & 89–97\% & 0.89–0.97 \\
This work — RF1                  & 4-class behavioural   & 97.9\%  & 0.988     \\
\hline
\end{tabular}
\end{table}

The honeypot classifier achieves a macro F1 of 0.986, which falls in
the upper range of published CICIDS2017 results despite addressing a
fundamentally different problem: CICIDS2017 classifiers operate on
network flow statistics computed over many packets, whereas the honeypot
classifier uses only six application-layer behavioural features measured
per session. This suggests that \textbf{timing and interaction behaviour
alone are sufficient to distinguish attacker profiles} with high
accuracy, without requiring network-level instrumentation.

It should be noted that the comparison is not direct: the honeypot
dataset uses ground-truth labels set by the generating bots, whereas
CICIDS2017 labels are analyst-assigned. The honeypot problem may
therefore be considered easier in terms of label quality, and harder in
terms of class granularity relative to binary anomaly detection tasks.

\section{Summary of Findings}

The experimental results demonstrate that:

\begin{enumerate}
    \item The multi-stage architecture successfully captures the complete
    attacker lifecycle across 38,438 sessions, with every phishing event
    linked to its corresponding SSH session via \texttt{session\_id}.

    \item Six behavioural features extracted from application-layer
    interactions are sufficient to classify attacker profiles with
    97.9\% accuracy and 98.8\% macro F1, meeting or exceeding
    CICIDS2017 benchmark performance.

    \item The four attacker profiles produce clearly separable
    behavioural fingerprints, with timing-based features
    (\texttt{typing\_delay\_ms}, \texttt{inter\_cmd\_delay\_avg\_s})
    contributing the most discriminating information.

    \item The multi-label architecture produces simultaneously useful
    intelligence: the attack stage classifier identifies \emph{where}
    in the lifecycle a session belongs, while the bot type classifier
    identifies \emph{who} is responsible.

    \item The macro F1 standard deviation of $\pm$0.1\% across
    cross-validation folds confirms that the results are stable and
    not an artefact of a particular train/test split.
\end{enumerate}
